\documentclass[10pt,a4paper]{article}
\usepackage[utf8]{inputenc}
\usepackage[margin=1.5cm]{geometry}
\usepackage{amsmath,amssymb,amsfonts}
\usepackage{array}
\usepackage{booktabs}
\usepackage{multicol}
\usepackage{tcolorbox}
\usepackage{xcolor}
\usepackage{hyperref}

% Konfigurasi warna estetis modern
\definecolor{primary}{RGB}{18, 52, 86}      % Navy Blue
\definecolor{secondary}{RGB}{28, 110, 164}  % Medium Blue
\definecolor{accent}{RGB}{217, 83, 79}      % Soft Coral
\definecolor{lightbg}{RGB}{245, 248, 252}   % Soft Blue-Gray
\definecolor{boxborder}{RGB}{180, 205, 230} % Border Blue
\definecolor{darktext}{RGB}{30, 41, 59}     % Slate Dark

\tcbuselibrary{skins,breakable}
\tcbset{
    enhanced,
    colback=lightbg,
    colframe=secondary,
    fonttitle=\bfseries\sffamily,
    coltitle=white,
    boxrule=0.8pt,
    arc=3mm,
    boxsep=2pt,
    left=5pt,right=5pt,top=4pt,bottom=4pt,
    before skip=6pt,
    after skip=6pt
}

\newtcolorbox{rumusbox}[2][]{
    colback=white,
    colframe=primary,
    title={#2},
    coltitle=white,
    fonttitle=\bfseries\sffamily\footnotesize,
    boxrule=0.8pt,
    arc=2mm,
    left=4pt,right=4pt,top=3pt,bottom=3pt,
    before skip=4pt,after skip=4pt,
    #1
}

\newtcolorbox{tipbox}[1][Tips \& Perangkap Ujian]{
    colback=white,
    colframe=accent,
    title={#1},
    coltitle=white,
    fonttitle=\bfseries\sffamily\footnotesize,
    boxrule=0.8pt,
    arc=2mm,
    left=4pt,right=4pt,top=3pt,bottom=3pt,
    before skip=4pt,after skip=4pt
}

\pagestyle{plain}

\begin{document}

% Header Utama
\begin{center}
    {\LARGE\bfseries\sffamily\color{primary} CHEATSHEET LENGKAP BEFORE UTS}\\[3pt]
    {\large\bfseries\sffamily\color{secondary} MA4171 Teori Kontrol Linear --- Semester I 2026/2027}\\[2pt]
    {\small Dosen: Dr. Dewi Handayani, S.Si., M.Si. \quad $\bullet$ \quad FMIPA Institut Teknologi Bandung}
\end{center}
\vspace{-6pt}
\hrule height 1.2pt
\vspace{6pt}

\begin{multicols*}{2}

% ========================================================
% BAGIAN 1: TRANSFORMASI LAPLACE & PDB
% ========================================================
\begin{rumusbox}{1. Transformasi Laplace \& PDB Linear}
\textbf{Definisi Dasar}:
$$\mathcal{L}\{f(t)\} = F(s) = \int_0^\infty f(t) e^{-st} \, dt$$

\textbf{Sifat-Sifat Krusial}:
\begin{itemize}\setlength\itemsep{1pt}
    \item \textbf{Turunan}: $\mathcal{L}\{\dot{x}(t)\} = s X(s) - x(0)$
    \item $\mathcal{L}\{\ddot{x}(t)\} = s^2 X(s) - s x(0) - \dot{x}(0)$
    \item $\mathcal{L}\{\dddot{x}(t)\} = s^3 X(s) - s^2 x(0) - s \dot{x}(0) - \ddot{x}(0)$
    \item $\mathcal{L}\{x^{(n)}\} = s^n X(s) - \sum_{k=1}^n s^{n-k} x^{(k-1)}(0)$
    \item \textbf{Integral}: $\mathcal{L}\left\{\int_0^t x(\tau) d\tau\right\} = \frac{1}{s} X(s)$
    \item \textbf{Pergeseran Waktu}: $\mathcal{L}\{f(t-a)u(t-a)\} = e^{-as} F(s)$
    \item \textbf{Pergeseran Frekuensi}: $\mathcal{L}\{e^{at} f(t)\} = F(s-a)$
    \item \textbf{Perkalian $t$}: $\mathcal{L}\{t f(t)\} = -\frac{d}{ds} F(s)$
\end{itemize}

\textbf{Tabel Pasangan Transformasi Penting}:
\begin{center}
\renewcommand{\arraystretch}{1.2}
\scriptsize
\begin{tabular}{cc|cc}
\toprule
$f(t)$ & $F(s)$ & $f(t)$ & $F(s)$ \\
\midrule
$\delta(t)$ & $1$ & $u(t)$ (step) & $\frac{1}{s}$ \\
$t$ & $\frac{1}{s^2}$ & $t^n$ & $\frac{n!}{s^{n+1}}$ \\
$e^{-at}$ & $\frac{1}{s+a}$ & $t e^{-at}$ & $\frac{1}{(s+a)^2}$ \\
$\sin(\omega t)$ & $\frac{\omega}{s^2+\omega^2}$ & $\cos(\omega t)$ & $\frac{s}{s^2+\omega^2}$ \\
$e^{-at}\sin(\omega t)$ & $\frac{\omega}{(s+a)^2+\omega^2}$ & $e^{-at}\cos(\omega t)$ & $\frac{s+a}{(s+a)^2+\omega^2}$ \\
\bottomrule
\end{tabular}
\end{center}

\textbf{Invers via Dekomposisi Pecahan Parsial}:
\begin{itemize}\setlength\itemsep{1pt}
    \item \textit{Akar Berbeda}: $\frac{P(s)}{(s-p_1)(s-p_2)} = \frac{A}{s-p_1} + \frac{B}{s-p_2}$
    $$A = \lim_{s \to p_1} (s-p_1) F(s)$$
    \item \textit{Akar Berulang}: $\frac{P(s)}{(s-p)^m} = \sum_{k=1}^m \frac{A_k}{(s-p)^k}$
    $$A_k = \frac{1}{(m-k)!} \lim_{s \to p} \frac{d^{m-k}}{ds^{m-k}} \left[ (s-p)^m F(s) \right]$$
\end{itemize}
\end{rumusbox}

% ========================================================
% BAGIAN 2: REDUKSI DIAGRAM BLOK
% ========================================================
\begin{rumusbox}{2. Aljabar Diagram Blok \& Fungsi Transfer}
\textbf{Fungsi Transfer Sistem SISO}:
$$G(s) = \frac{Y(s)}{U(s)} \quad \text{dengan kondisi awal } x(0) = 0$$

\textbf{Tiga Hubungan Dasar}:
\begin{itemize}\setlength\itemsep{2pt}
    \item \textbf{Kaskade (Seri)}: $G_{eq}(s) = G_1(s) G_2(s)$
    \item \textbf{Paralel}: $G_{eq}(s) = G_1(s) \pm G_2(s)$
    \item \textbf{Umpan Balik (\textit{Feedback})}:
    $$G_{eq}(s) = \frac{C(s)}{R(s)} = \frac{G(s)}{1 \pm G(s)H(s)}$$
    $\bullet$ Umpan balik negatif (tanda $-$ di pembanding): $+G(s)H(s)$\\
    $\bullet$ Umpan balik positif (tanda $+$ di pembanding): $-G(s)H(s)$
\end{itemize}

\textbf{Aturan Ekuivalensi Pemindahan}:
\begin{itemize}\setlength\itemsep{1pt}
    \item Pindah titik cabang \textbf{melewati blok $G$ ke depan}: kalikan cabang dengan $\frac{1}{G}$.
    \item Pindah titik cabang \textbf{melewati blok $G$ ke belakang}: kalikan cabang dengan $G$.
    \item Pindah titik penjumlahan \textbf{melewati blok $G$ ke depan}: kalikan masukan baru dengan $G$.
    \item Pindah titik penjumlahan \textbf{melewati blok $G$ ke belakang}: kalikan masukan baru dengan $\frac{1}{G}$.
\end{itemize}
\end{rumusbox}

% ========================================================
% BAGIAN 3: ATURAN PENGUATAN MASON (SFG)
% ========================================================
\begin{rumusbox}{3. Signal Flow Graph \& Aturan Mason}
\textbf{Rumus Penguatan Mason (\textit{Mason's Gain Formula})}:
$$P = \frac{Y_{out}}{U_{in}} = \frac{1}{\Delta} \sum_{k=1}^N P_k \Delta_k$$

\textbf{Komponen Formula}:
\begin{itemize}\setlength\itemsep{2pt}
    \item $P_k$ : Penguatan lintasan maju (\textit{forward path gain}) ke-$k$ dari node input ke node output tanpa menyentuh node mana pun dua kali.
    \item $\Delta$ : Determinan graf keseluruhan:
    $$\Delta = 1 - \sum L_i + \sum L_i L_j - \sum L_i L_j L_k + \dots$$
    di mana:
    \begin{itemize}\setlength\itemsep{1pt}
        \item $\sum L_i$ : jumlah seluruh gain loop individu.
        \item $\sum L_i L_j$ : jumlah perkalian gain pasangan dua loop yang \textbf{tidak saling bersentuhan} (\textit{non-touching}).
        \item $\sum L_i L_j L_k$ : jumlah perkalian 3 loop yang saling tidak bersentuhan.
    \end{itemize}
    \item $\Delta_k$ : Kofaktor lintasan ke-$k$, yaitu nilai determinan graf $\Delta$ setelah \textbf{menghapus seluruh loop yang bersentuhan} dengan lintasan maju $P_k$.
\end{itemize}
\end{rumusbox}

\begin{tipbox}[Prosedur Cepat Menjawab Soal Mason di Ujian]
1. Tuliskan semua lintasan maju $P_1, P_2, \dots$ beserta nilainya.\\
2. Tuliskan seluruh loop tunggal $L_1, L_2, \dots$ dengan tanda aljabar yang tepat (perhatikan tanda $-H$).\\
3. Identifikasi loop yang \textit{non-touching} (tidak berbagi simpul).\\
4. Hitung $\Delta$, lalu tentukan $\Delta_k$ (jika lintasan $P_k$ menyentuh semua loop, maka $\Delta_k = 1$).\\
5. Masukkan ke rumus $P = \frac{\sum P_k \Delta_k}{\Delta}$.
\end{tipbox}

% ========================================================
% BAGIAN 4: RUANG KEADAAN & RESPON WAKTU
% ========================================================
\begin{rumusbox}{4. Representasi Ruang Keadaan \& Respon Sistem}
\textbf{Bentuk Standar LTI}:
\begin{align*}
    \dot{x}(t) &= A x(t) + B u(t) \quad (\text{Persamaan Keadaan})\\
    y(t) &= C x(t) + D u(t) \quad (\text{Persamaan Keluaran})
\end{align*}
dengan $x \in \mathbb{R}^n, u \in \mathbb{R}^m, y \in \mathbb{R}^p$.

\textbf{Konversi Ruang Keadaan ke Fungsi Transfer}:
$$G(s) = C(sI - A)^{-1} B + D = \frac{C \operatorname{adj}(sI - A) B}{\det(sI - A)} + D$$
Kutub-kutub (\textit{poles}) sistem adalah akar-akar dari $\det(sI - A) = 0$.

\textbf{Matriks Fundamental / Transisi Keadaan $\Phi(t) = e^{At}$}:
$$e^{At} = \mathcal{L}^{-1}\{(sI - A)^{-1}\}$$
Untuk matriks $2 \times 2$:
$$(sI - A)^{-1} = \frac{1}{\det(sI - A)} \begin{bmatrix} s - a_{22} & a_{12} \\ a_{21} & s - a_{11} \end{bmatrix}$$

\textbf{Solusi Lengkap Respon Waktu}:
$$x(t) = \underbrace{e^{At} x(0)}_{\text{Respon Masukan Nol (ZIR)}} + \underbrace{\int_0^t e^{A(t-\tau)} B u(\tau) \, d\tau}_{\text{Respon Keadaan Nol (ZSR)}}$$
Dalam domain Laplace:
$$X(s) = (sI - A)^{-1} x(0) + (sI - A)^{-1} B U(s)$$
$$Y(s) = C(sI - A)^{-1} x(0) + \left[ C(sI - A)^{-1} B + D \right] U(s)$$
\end{rumusbox}

% ========================================================
% BAGIAN 5: BENTUK KANONIK RUANG KEADAAN
% ========================================================
\begin{rumusbox}{5. Empat Bentuk Kanonik Ruang Keadaan}
Diberikan fungsi transfer SISO berordo $n$:
$$G(s) = \frac{b_0 s^n + b_1 s^{n-1} + \dots + b_{n-1}s + b_n}{s^n + a_1 s^{n-1} + \dots + a_{n-1}s + a_n}$$
Bila $b_0 = 0 \implies D = 0$.

\textbf{1. Bentuk Kanonik Keterkontrolan (CCF)}:\\
Variabel perantara: $Z(s) = \frac{U(s)}{s^n + a_1 s^{n-1} + \dots + a_n}$.\\
State: $x_1 = z, x_2 = \dot{z}, \dots, x_n = z^{(n-1)}$.
{\scriptsize
\begin{align*}
A &= \begin{bmatrix} 0 & 1 & \dots & 0 \\ \vdots & \vdots & \ddots & \vdots \\ 0 & 0 & \dots & 1 \\ -a_n & -a_{n-1} & \dots & -a_1 \end{bmatrix}, \quad B = \begin{bmatrix} 0 \\ \vdots \\ 0 \\ 1 \end{bmatrix} \\
C &= \begin{bmatrix} b_n - a_n b_0 & \dots & b_1 - a_1 b_0 \end{bmatrix}, \quad D = b_0
\end{align*}
}

\textbf{2. Bentuk Kanonik Keterobservasian (OCF)}:\\
Bentuk dual dari CCF ($A_{ocf} = A_{ccf}^T, B_{ocf} = C_{ccf}^T, C_{ocf} = B_{ccf}^T$):
{\scriptsize
\begin{align*}
A &= \begin{bmatrix} 0 & \dots & 0 & -a_n \\ 1 & \dots & 0 & -a_{n-1} \\ \vdots & \ddots & \vdots & \vdots \\ 0 & \dots & 1 & -a_1 \end{bmatrix}, \quad B = \begin{bmatrix} b_n - a_n b_0 \\ b_{n-1} - a_{n-1} b_0 \\ \vdots \\ b_1 - a_1 b_0 \end{bmatrix} \\
C &= \begin{bmatrix} 0 & \dots & 0 & 1 \end{bmatrix}, \quad D = b_0
\end{align*}
}

\textbf{3. Bentuk Kanonik Diagonal (DCF)}:\\
Syarat: Nilai eigen $\lambda_1, \dots, \lambda_n$ bernilai riil dan berbeda.\\
Ekspansi parsial: $G(s) = \sum_{i=1}^n \frac{c_i}{s - \lambda_i} + D$.
{\scriptsize
$$A = \operatorname{diag}(\lambda_1, \dots, \lambda_n), \quad B = \begin{bmatrix} 1 \\ \vdots \\ 1 \end{bmatrix}, \quad C = \begin{bmatrix} c_1 & \dots & c_n \end{bmatrix}$$
}

\textbf{4. Bentuk Kanonik Jordan (JCF)}:\\
Bila ada nilai eigen berulang, misal $\lambda_1$ rangkap 3:
{\scriptsize
$$J_1 = \begin{bmatrix} \lambda_1 & 1 & 0 \\ 0 & \lambda_1 & 1 \\ 0 & 0 & \lambda_1 \end{bmatrix}, \quad B = \begin{bmatrix} 0 \\ 0 \\ 1 \end{bmatrix}$$
}
\end{rumusbox}

% ========================================================
% BAGIAN 6: SISTEM MULTIVARIABEL (SIMO, MISO, MIMO)
% ========================================================
\begin{rumusbox}{6. Sistem Multivariabel (SIMO, MISO, MIMO)}
\begin{itemize}\setlength\itemsep{1pt}
    \item \textbf{SIMO (1 masukan, $p$ keluaran)}:
    $$G(s) = \frac{1}{d(s)} \begin{bmatrix} n_1(s) \\ \vdots \\ n_p(s) \end{bmatrix}$$
    Gunakan CCF bersama: matriks $A$ dan $B$ sama, $C \in \mathbb{R}^{p \times n}$ memuat koefisien masing-masing $n_i(s)$.
    \item \textbf{MISO ($m$ masukan, 1 keluaran)}:
    $$G(s) = \frac{1}{d(s)} \begin{bmatrix} n_1(s) & \dots & n_m(s) \end{bmatrix}$$
    Gunakan OCF bersama: matriks $A$ dan $C$ sama, $B \in \mathbb{R}^{n \times m}$ memuat koefisien masing-masing $n_i(s)$.
    \item \textbf{MIMO ($m$ masukan, $p$ keluaran)}:
    Matriks fungsi transfer $\mathbf{G}(s) \in \mathbb{R}^{p \times m}$. Realisasi minimal dapat dibentuk melalui bentuk Diagonal/Jordan via residu parsial.
\end{itemize}
\end{rumusbox}

% ========================================================
% BAGIAN 7: KETERKONTROLAN & KETEROBSERVASIAN (TOPIK 04)
% ========================================================
\begin{rumusbox}{7. Keterkontrolan \& Keterobservasian Sistem}
\textbf{1. Kriteria Keterkontrolan Kalman}:
Matriks keterkontrolan:
$$M_c = \begin{bmatrix} B & AB & A^2B & \dots & A^{n-1}B \end{bmatrix} \in \mathbb{R}^{n \times n}$$
Sistem terkontrol lengkap $\iff \operatorname{rank}(M_c) = n$ (atau $\det(M_c) \ne 0$).

\textbf{2. Kriteria Keterobservasian Kalman}:
Matriks keterobservasian:
$$M_o = \begin{bmatrix} C \\ CA \\ CA^2 \\ \vdots \\ CA^{n-1} \end{bmatrix} \in \mathbb{R}^{n \times n}$$
Sistem terobservasi lengkap $\iff \operatorname{rank}(M_o) = n$ (atau $\det(M_o) \ne 0$).

\textbf{3. Transformasi ke CCF Menggunakan $T = M_c W$}:
Jika $|sI - A| = s^n + a_1 s^{n-1} + \dots + a_n = 0$:
{\scriptsize
$$W = \begin{bmatrix} a_{n-1} & a_{n-2} & \dots & 1 \\ a_{n-2} & \dots & 1 & 0 \\ \vdots & \reflectbox{$\ddots$} & \vdots & \vdots \\ 1 & 0 & \dots & 0 \end{bmatrix}, \quad T = M_c W$$
}
Transformasi $x = Tz$ menghasilkan:
$$A_{ccf} = T^{-1} A T, \quad B_{ccf} = T^{-1} B$$

\textbf{4. Uji Gilbert (Bentuk Diagonal \& Jordan)}:
\begin{itemize}\setlength\itemsep{1pt}
    \item \textbf{Bentuk Diagonal} ($A = \operatorname{diag}(\lambda_i)$):
    \begin{itemize}\setlength\itemsep{0pt}
        \item Terkontrol $\iff$ Tidak ada baris pada $\tilde{B}$ yang seluruhnya bernilai nol.
        \item Terobservasi $\iff$ Tidak ada kolom pada $\tilde{C}$ yang seluruhnya bernilai nol.
    \end{itemize}
    \item \textbf{Bentuk Jordan} (Blok $J_i$ nilai eigen berulang):
    \begin{itemize}\setlength\itemsep{0pt}
        \item Terkontrol $\iff$ Elemen pada \textbf{baris terakhir} setiap blok Jordan pada matriks $\tilde{B}$ bernilai bukan nol ($b_{terakhir} \ne 0$).
        \item Terobservasi $\iff$ Elemen pada \textbf{kolom pertama} setiap blok Jordan pada matriks $\tilde{C}$ bernilai bukan nol ($c_{pertama} \ne 0$).
    \end{itemize}
\end{itemize}

\textbf{5. Prinsip Dualitas Kalman}:
$$(A, B) \text{ terkontrol} \iff (A^T, B^T) \text{ terobservasi}$$
$$M_c(A, B) = M_o^T(A^T, B^T)$$
\end{rumusbox}

% ========================================================
% BAGIAN 8: TRANSFORMASI Z & SISTEM WAKTU DISKRIT
% ========================================================
\begin{rumusbox}{8. Transformasi Z \& Sistem Waktu Diskrit}
\textbf{Definisi}:
$$X(z) = \mathcal{Z}\{x(k)\} = \sum_{k=0}^\infty x(k) z^{-k}$$

\textbf{Tabel Transformasi Z Esensial}:
\begin{center}
\renewcommand{\arraystretch}{1.2}
\scriptsize
\begin{tabular}{cc|cc}
\toprule
$x(k)$ & $X(z)$ & $x(k)$ & $X(z)$ \\
\midrule
$\delta(k)$ & $1$ & $1$ (step) & $\frac{z}{z-1} = \frac{1}{1-z^{-1}}$ \\
$a^k$ & $\frac{z}{z-a}$ & $k$ & $\frac{z}{(z-1)^2}$ \\
$k a^k$ & $\frac{a z}{(z-a)^2}$ & $a^{k-1}$ & $\frac{1}{z-a}$ \\
\bottomrule
\end{tabular}
\end{center}

\textbf{Teorema Pergeseran Waktu}:
\begin{align*}
    \mathcal{Z}\{x(k+1)\} &= z X(z) - z x(0)\\
    \mathcal{Z}\{x(k+2)\} &= z^2 X(z) - z^2 x(0) - z x(1)
\end{align*}

\textbf{Metode Inversi Pecahan Parsial}:
Uraikan selalu dalam bentuk $\frac{X(z)}{z}$:
$$\frac{X(z)}{z} = \frac{A}{z - p_1} + \frac{B}{z - p_2} \implies X(z) = \frac{A z}{z - p_1} + \frac{B z}{z - p_2}$$
sehingga dengan mudah diperoleh $x(k) = A p_1^k + B p_2^k$.

\textbf{Ruang Keadaan Diskrit}:
\begin{align*}
    x(k+1) &= G x(k) + H u(k)\\
    y(k) &= C x(k) + D u(k)
\end{align*}
Fungsi transfer diskrit: $G(z) = C(zI - G)^{-1}H + D$.
\end{rumusbox}

\begin{tipbox}[Checklist Final Peserta Ujian UTS]
$\square$ Selalu tulis asumsi kondisi awal nol saat mencari fungsi transfer.\\
$\square$ Jangan tertukar: Umpan balik negatif bertanda $+$ di penyebut $\frac{G}{1+GH}$.\\
$\square$ Saat menghitung determinan $\Delta$ Mason, periksa dengan teliti apakah ada loop yang saling lepas (\textit{non-touching}).\\
$\square$ Untuk matriks $e^{At}$, pastikan invers Laplace $(sI - A)^{-1}$ dihitung elemen demi elemen dengan benar.\\
$\square$ Rank matriks $M_c$ atau $M_o$: jika determinan $\ne 0$, rank langsung $n$ (penuh).
\end{tipbox}

\end{multicols*}

\end{document}
